\documentclass{ctexart}
\usepackage[paperwidth=235mm,paperheight=78mm,margin=4mm]{geometry}
\usepackage{tikz}
\pagestyle{empty}
\begin{document}
\thispagestyle{empty}
\begin{tikzpicture}[x=0.95cm,y=0.95cm]
\usetikzlibrary{arrows.meta,positioning,calc,fit,backgrounds,shapes.geometric,shapes.symbols}

\tikzset{
  chfive/box/.style={
    draw,
    rounded corners=3pt,
    thick,
    align=center,
    minimum height=8mm,
    minimum width=24mm,
    font=\small,
    fill=blue!8
  },
  chfive/core/.style={
    draw,
    rounded corners=3pt,
    thick,
    align=center,
    minimum height=8mm,
    minimum width=26mm,
    font=\small\bfseries,
    fill=orange!15
  },
  chfive/aux/.style={
    draw,
    rounded corners=3pt,
    thick,
    align=center,
    minimum height=8mm,
    minimum width=24mm,
    font=\small,
    fill=green!10
  },
  chfive/data/.style={
    draw,
    rounded corners=3pt,
    thick,
    align=center,
    minimum height=8mm,
    minimum width=24mm,
    font=\small,
    fill=yellow!16
  },
  chfive/group/.style={
    draw,
    rounded corners=4pt,
    thick,
    dashed,
    inner sep=6pt
  },
  chfive/arr/.style={-{Latex[length=2.6mm]}, thick},
  chfive/darr/.style={-{Latex[length=2.6mm]}, thick, dashed},
  chfive/lab/.style={font=\small},
  chfive/tile/.style={draw=black!60, line width=0.5pt, fill=gray!5},
  chfive/corridor/.style={draw=black!70, line width=1.0pt, fill=cyan!12},
  chfive/highlight/.style={draw=red!75!black, line width=1.0pt, fill=red!12},
  chfive/active/.style={draw=orange!80!black, line width=1.0pt, fill=orange!20},
  chfive/traj/.style={draw=blue!70!black, line width=1.2pt, -{Latex[length=2.4mm]}},
  chfive/boundary/.style={draw=red!70!black, dashed, line width=1pt},
  chfive/chartA/.style={draw=blue!70!black, fill=blue!30},
  chfive/chartB/.style={draw=orange!80!black, fill=orange!30}
}


% chart 1
\begin{scope}[shift={(0,0)}]
  \draw[->, thick] (0,0) -- (4.4,0) node[below] {$R$};
  \draw[->, thick] (0,0) -- (0,3.4) node[left] {归一化指标};
  \fill[chfive/chartA] (0.6,0) rectangle (1.1,1.0);
  \fill[chfive/chartA] (1.7,0) rectangle (2.2,1.6);
  \fill[chfive/chartA] (2.8,0) rectangle (3.3,2.3);
  \node[chfive/lab] at (2.0,-0.7) {内存/计算开销随 $R$ 增大};
  \foreach \x/\lab in {0.85/1,1.95/2,3.05/3} \node[chfive/lab, below] at (\x,0) {\lab};
\end{scope}

% chart 2
\begin{scope}[shift={(6.3,0)}]
  \draw[->, thick] (0,0) -- (4.4,0) node[below] {$d_{\mathrm{th}}$};
  \draw[->, thick] (0,0) -- (0,3.4) node[left] {切换频率};
  \draw[chfive/chartB, thick] (0.5,2.8) -- (1.4,2.0) -- (2.3,1.3) -- (3.3,0.9);
  \foreach \p in {(0.5,2.8),(1.4,2.0),(2.3,1.3),(3.3,0.9)} {\fill[orange!80!black] \p circle (1.8pt);} 
  \node[chfive/lab] at (2.0,-0.7) {切换频率随门控阈值降低};
\end{scope}

% chart 3
\begin{scope}[shift={(12.5,0)}]
  \draw[->, thick] (0,0) -- (4.4,0) node[below] {$N_h$};
  \draw[->, thick] (0,0) -- (0,3.4) node[left] {稳定性};
  \draw[chfive/chartA, thick] (0.5,0.8) -- (1.4,1.5) -- (2.3,2.1) -- (3.3,2.5);
  \foreach \p in {(0.5,0.8),(1.4,1.5),(2.3,2.1),(3.3,2.5)} {\fill[blue!70!black] \p circle (1.8pt);} 
  \node[chfive/lab] at (2.0,-0.7) {滞回确认增强切换稳定性};
\end{scope}
\end{tikzpicture}
\end{document}
